\documentclass{report}
\usepackage{mathesis}
\usepackage[utf8]{inputenc}
\usepackage[russian]{babel}

\begin{document}

\section{Предел последовательности (Limit of Sequence)}\label{entity:op-limit-seq}
\textbf{Тип:} operation\quad \textbf{Источник:} Зорич В.А., Математический анализ, Том I (zorich-1, p. 77)

\textbf{Каноническая запись:}

\textbf{Естественный язык:}
Число $L \in \mathbb{R}$ называется пределом последовательности $(x_n)_{n \in \mathbb{N}}$, если для любого, даже сколь угодно малого, $\varepsilon > 0$ существует такой номер $n_0 \in \mathbb{N}$, что для всех членов последовательности с номерами $n \geq n_0$ выполняется неравенство $\entityref{op-abs-abstract}{\mathrm{abs}}(x_n - L) < \varepsilon$. 

\begin{operation}[op-limit-seq]
$(x_n)_{n \in \mathbb{N}} \in \mathbb{R}^{\mathbb{N}}, L \in \mathbb{R}$

\[
\mLim_{n \to \infty} x_n = L \mDefIff \mForall{\varepsilon > 0}\; \mExists{n_0 \in \mathbb{N}}\; \mForall{n \in \mathbb{N}} \colon \left( n \geq n_0 \mImplies \entityref{op-abs-abstract}{\mathrm{abs}}(x_n - L) < \varepsilon \right)
\]
\end{operation}


\vspace{1em}
\hrule
\vspace{1em}

\section{Абсолютная величина (Absolute Value)}\label{entity:op-abs-abstract}
\textbf{Тип:} operation\quad \textbf{Источник:} Зорич В.А., Математический анализ, Том I (zorich-1, p. 1)

\textbf{Каноническая запись:}

\begin{operation}[op-abs-abstract]

$x \in \entityref{obj-real-numbers}{\mathbb{R}}$

\[ |x| \mDefIff \max(x, -x) \]
\end{operation}


\vspace{1em}
\hrule
\vspace{1em}

\section{Вещественные числа (Real Numbers)}\label{entity:obj-real-numbers}
\textbf{Тип:} object\quad \textbf{Источник:} Зорич В.А., Математический анализ, Том I (zorich-1, p. 1)

\textbf{Каноническая запись:}

\begin{object}[obj-real-numbers]



\[ \mathbb{R} = (\entityref{obj-set}{R}, +, \cdot, \leq) \]
\end{object}


\textbf{Естественный язык:}
Множество вещественных чисел $\mathbb{R}$ — это непрерывная числовая прямая. Формально $\mathbb{R}$ задается как полное (непрерывное) архимедово упорядоченное поле. В нем можно складывать, умножать, сравнивать элементы, и в нем нет «дырок» (каждое ограниченное сверху подмножество имеет точную верхнюю грань).

\vspace{1em}
\hrule
\vspace{1em}

\section{Множество (Set)}\label{entity:obj-set}
\textbf{Тип:} object\quad \textbf{Источник:} Зорич В.А., Математический анализ, Том I (zorich-1, p. 7)

\textbf{Каноническая запись:}




\textbf{Естественный язык:}
Множество — базовое, неопределяемое напрямую понятие математики. Множество представляет собой совокупность объектов произвольной природы, называемых его элементами. Все свойства множеств строго выводятся из аксиом системы Цермело-Френкеля (ZFC).

\vspace{1em}
\hrule
\vspace{1em}

\section{Аксиома объемности (Axiom of Extensionality)}\label{entity:axm-zfc-extensionality}
\textbf{Тип:} axiom\quad \textbf{Источник:} vereshchagin-shen (vereshchagin-shen)

\textbf{Каноническая запись:}

\[
\mForall{A, B}
(A = B) \mIff (\mForall{x} x \in A \mIff x \in B)
\]


\vspace{1em}
\hrule
\vspace{1em}

\section{Аксиома объединения (Axiom of Union)}\label{entity:axm-zfc-union}
\textbf{Тип:} axiom\quad \textbf{Источник:} vereshchagin-shen (vereshchagin-shen)

\textbf{Каноническая запись:}

\[
\mForall{\mathcal{F}}
\mExists{U} \mForall{x} \left( x \in U \mIff \mExists{A} (A \in \mathcal{F} \entityref{def-logical-connectives}{\land} x \in A) \right)
\]


\vspace{1em}
\hrule
\vspace{1em}

\section{Аксиома пары (Axiom of Pairing)}\label{entity:axm-zfc-pairing}
\textbf{Тип:} axiom\quad \textbf{Источник:} vereshchagin-shen (vereshchagin-shen)

\textbf{Каноническая запись:}

\[
\mForall{A, B}
\mExists{C} \mForall{x} \left( x \in C \mIff (x = A \entityref{def-logical-connectives}{\lor} x = B) \right)
\]


\vspace{1em}
\hrule
\vspace{1em}



\end{document}
